
\subsubsection{6.1 Interpretation}

\textbf{The bias is large, structural, and persistent.} Cramér's V = 0.40--0.57 places the language $\times$ retention association firmly in the ''large'' effect range. The consistent es > fr > it > en > de ordering across all five threshold levels is not compatible with a sampling artifact; it tracks the Romance--Germanic language family boundary precisely. At k=30, a practitioner would retain 6.$3\times$ more Spanish documents than German documents from the same raw corpus. The resulting training corpus does not reflect relative document quality; it reflects a calibration artifact introduced by applying a global threshold to locally-calibrated scores.

\textbf{Prior estimates underestimate the bias.} The finding that V = 0.40--0.57 exceeds literature-derived estimates by 25--40\% has direct consequences for correction study design. Researchers using the CCNet paper or analogous descriptions to size their correction experiments will underestimate the required effect magnitude. The calibrated baseline V = 0.40--0.57 on RedPajama-V2 ccnet\_perplexity should be used when designing correction experiments on this signal.

\textbf{Connection to prior work.} This measurement extends Caswell et al. [2021]'s qualitative documentation of multilingual corpus quality disparities to a specific quantitative baseline on ccnet\_perplexity. It confirms the practical consequence of deviating from CCNet's original per-language design intent \citep{Wenzek2019CCNet} by precisely quantifying the resulting disparity. It provides the V baseline that Turki et al. [2026]'s endorsement of percentile-based per-language thresholding implicitly requires for rigorous evaluation.

\subsubsection{6.2 Limitations}

\textbf{L1: Existence-only scope.} This paper tests only the existence and magnitude of the disparity under global thresholding. Whether per-language calibration reduces the disparity substantially ($\Delta V \geq$ 0.10 for $\geq 3/5$ k values) remains to be tested in follow-on work. The characterization contribution is complete and independently publishable as a necessary baseline.

\textbf{L2: Sample scope.} Results are based on 208,262 documents, approximately 0.0002\% of the full 113.3 billion document RedPajama-V2 corpus. The phenomenon is structurally driven and confirmed across three independent runs on the same sample; extrapolation to the full corpus requires distributed-scale replication.

\textbf{L3: Document-length confound.} Longer documents tend to have lower perplexity. Length stratification was planned but not implemented in the current analysis. The 72.7 percentage-point max--min retention gap substantially exceeds what a plausible length confound alone can explain; however, length stratification is a recommended robustness check for follow-on work.

\textbf{L4: Five European languages.} The five-language sample covers only the Germanic (en, de) and Romance (es, fr, it) families within Indo-European. Findings may not generalize to low-resource, agglutinative, or non-Indo-European languages, which may exhibit different or more severe distribution divergences from Wikipedia-trained KenLM models.

\subsubsection{6.3 Broader Impact}

Any model trained on a corpus filtered using global percentile thresholds on ccnet\_perplexity will have structurally imbalanced language coverage, favoring Romance-family languages over Germanic-family languages regardless of the specific threshold chosen. This work enables practitioners to: (a) audit existing multilingual training datasets for this bias source, (b) design properly powered correction studies using the V = 0.40--0.57 calibrated baseline, and (c) adopt per-language calibration as a default practice. Correction strategies should be evaluated against the baseline established here before deployment in production pipelines.

